% ID: SEL-D0073-OB-E02
% Tanda S073.
\textbf{a)} La recta tiene pendiente $(0-(-3))/(4-0)=3/4$ y ordenada en el origen $-3$, de modo que $f'(x)=3x/4-3$. Es negativa para $x<4$ y positiva para $x>4$. Por el criterio del signo de la derivada, $f$ decrece en $(-\infty,4)$ y crece en $(4,\infty)$. Tiene un mínimo en $x=4$.

\textbf{b)} Interpretamos $L$ como logaritmo natural. Aplicamos producto y cadena a $g$, y cociente a $h$:
\[
 \begin{aligned}
 g'(x)&=9(3x+1)^2\ln(x^2+1)+(3x+1)^3\frac{2x}{x^2+1},\\
 h'(x)&=\frac{e^x(7x^5-4)-e^x\cdot35x^4}{(7x^5-4)^2}
      =\frac{e^x(7x^5-35x^4-4)}{(7x^5-4)^2}.
 \end{aligned}
\]
La primera fórmula vale para todo $x$ real y la segunda para $x\ne\sqrt[5]{4/7}$.
